\documentclass{article}
\usepackage{amsmath,amssymb}
\usepackage{xurl}
\usepackage[colorlinks=true,linkcolor=blue,citecolor=blue,urlcolor=blue]{hyperref}
% Shared commands for notes in docs/paper-gaps/.

\DeclareUnicodeCharacter{2080}{\ifmmode{}_{0}\else\textsubscript{0}\fi}
\DeclareUnicodeCharacter{2081}{\ifmmode{}_{1}\else\textsubscript{1}\fi}
\DeclareUnicodeCharacter{2082}{\ifmmode{}_{2}\else\textsubscript{2}\fi}
\DeclareUnicodeCharacter{2099}{\ifmmode{}_{n}\else\textsubscript{n}\fi}

\newcommand{\Lean}{\textsc{Lean}}
\ExplSyntaxOn
\NewDocumentCommand{\leanid}{m}{
  \ifmmode
    \textsf{\detokenize{#1}}
  \else
    \group_begin:
    \urlstyle{sf}
    \tl_set:Nn \l_tmpa_tl {#1}
    \tl_replace_all:Nnn \l_tmpa_tl {\_} {_}
    \exp_args:NV \path \l_tmpa_tl
    \group_end:
  \fi
}
\ExplSyntaxOff
\providecommand{\tr}{\operatorname{tr}}
\newcommand{\tnleanIssueBase}{https://github.com/LionSR/TNLean/issues/}
\newcommand{\tnleanPrBase}{https://github.com/LionSR/TNLean/pull/}
\newcommand{\tnleanDocBase}{https://sirui-lu.com/TNLean/docs/}
\newcommand{\ghissue}[1]{\href{\tnleanIssueBase#1}{GitHub issue \##1}}
\newcommand{\ghpr}[1]{\href{\tnleanPrBase#1}{PR \##1}}
\newcommand{\doclink}[1]{\href{\tnleanDocBase#1.html}{\path{#1}}}

% Dirac notation and the identity operator, as in blueprint/src/macros/common.tex.
% \providecommand keeps note-local definitions authoritative.
\usepackage{dsfont}
\providecommand{\ket}[1]{|#1\rangle}
\providecommand{\bra}[1]{\langle #1|}
\providecommand{\braket}[2]{\langle #1 | #2 \rangle}
\providecommand{\ketbra}[2]{|#1\rangle\!\langle #2|}
\providecommand{\Id}{\mathds{1}}
\providecommand{\one}{\mathds{1}}
\providecommand{\C}{\mathbb{C}}
\providecommand{\MN}[1]{M_{#1}(\C)}
\providecommand{\MD}{M_D(\C)}

% Verdict marker: \gapnote{<kind>}{<status>}. Kind is the policy
% classification (clarification, local-correction, scope-restriction,
% unfaithful, false-source, open-gap); status is open, wip (elimination
% underway), resolved, or historical. Machine-read by the site index;
% prints a verdict line.
\newcommand{\gapnote}[2]{\par\noindent\textbf{Verdict:} #1 (#2).\par}

\setlength{\emergencystretch}{2em}
\title{The Decay Length of Lemma~1 in Hastings' Area Law}
\author{}
\date{October 28, 2026}
\begin{document}
\maketitle
\gapnote{local-correction}{resolved}

\section{The source statement}

The main theorem of~\cite{Hastings2007AreaLaw} and its Lemma~1 use the
lengths
\begin{align}
    \xi=\max(2v/\Delta E,\xi_C),\qquad\xi'=6\xi,
    \label{eq:hastings07_lemma1_decay_length_source}
\end{align}
where $\Delta E$ is the gap and $v$, $\xi_C$ are a velocity and a length
scale of a Lieb--Robinson bound.\footnote{\path{References/0705.2024/law16.tex},
    lines~139--142.} That bound is printed as
\begin{align}
    \|[A(t),B]\|\le c\,|X|\,\|A\|\,\|B\|\,e^{-\xi_C\operatorname{dist}(X,Y)}
    \quad\text{for }|t|\le l/v,
    \label{eq:hastings07_lemma1_decay_length_lrb}
\end{align}
for operators $A$, $B$ supported on sets $X$, $Y$; the source says only that
such $v$ and $\xi_C$ exist, with $v$ of order $J$ and $\xi_C$ of order
one.\footnote{\path{References/0705.2024/law16.tex}, lines~96--112.}

\section{The reading}

As printed, (\ref{eq:hastings07_lemma1_decay_length_lrb}) multiplies the
distance by $\xi_C$ and bounds the time by $l/v$, a length not defined at
that point. The reading consistent with the name ``length scale'' and with
the use of $2v/\Delta E$ as a competing length is that the commutator decays
as $e^{-\operatorname{dist}(X,Y)/\xi_C}$ for $|t|\le\operatorname{dist}(X,Y)/v$.
Any pair $(v,\xi_C)$ for which this holds may be used.

\section{The numerical chain estimate}

For range-$R$ interactions of norm at most $J\ge0$, put
$c_R=2(R+1)2^{2R+1}e^{2R}$ and $v_0=1+c_RJ$.
The numerical finite-support propagation estimate gives
\begin{align}
    \|[A(t),B]\|\le4|X|\,\|A\|\,\|B\|
        e^{2c_RJ|t|-d},
    \label{eq:hastings07_lemma1_decay_length_chain}
\end{align}
whenever every distance between a site of $X$ and a site of $Y$ is at least
$d$. Both supports are arbitrary finite sets; there is no factor $|Y|$.
The prefactor is independent of the interaction strength, local dimension
and chain length.\footnote{\leanid{TNLean.MPS.AreaLaw.ChainHamiltonian.norm_heisenberg_commutator_le_numeric}
    in \doclink{TNLean/MPS/AreaLaw/Hastings/LiebRobinson}.}
If $v>0$, $v\ge8c_RJ$, $\xi_C\ge8$ and $|t|\le d/v$, then
$2c_RJ|t|\le d/4$ and $d/\xi_C\le d/8$. Thus
\begin{align}
    \|[A(t),B]\|\le4|X|\,\|A\|\,\|B\|e^{-d/\xi_C}.
    \label{eq:hastings07_lemma1_decay_length_cone}
\end{align}
This proves the source's time-cone bound in the reciprocal-length
convention, including monotonicity under increasing the positive velocity
or the length scale.\footnote{\leanid{TNLean.MPS.AreaLaw.ChainHamiltonian.norm_heisenberg_commutator_le_hastings}.}
The distance lower-bound formulation includes empty supports without
assigning a finite distance to the empty set.

The localization proof used to construct the factors still uses its
conditional-expectation estimate. That estimate carries $|X|^2$, but the
support has at most $6R$ sites for the filtered parts. For the boundary
operator its size is $O(R+l)$, and the faster decay of that term absorbs
the resulting polynomial in $l$. This affects the constant rather than the
admissibility of the velocity and length pair.

\section{Energy units and the two factorization lengths}

The unscaled construction uses
\begin{align}
    \xi=8\max(2v_0/\Delta E,1)=\max(16v_0/\Delta E,8).
    \label{eq:hastings07_lemma1_decay_length_value}
\end{align}
Its velocity and length pair is $(8v_0,8)$, which satisfies
(\ref{eq:hastings07_lemma1_decay_length_cone}).
For the construction in units of the gap, let
$u_0=1+c_RJ/\Delta E$. Then
\begin{align}
    \hat\xi=16u_0=16(1+c_RJ/\Delta E).
    \label{eq:hastings07_lemma1_decay_length_scaled}
\end{align}
The physical velocity is $8\Delta E\,u_0$, so
$\hat\xi=\max(2(8\Delta E\,u_0)/\Delta E,8)$.
The factor $\Delta E$ is necessary: evolution under $H/\Delta E$ at time
$\tau=\Delta E\,t$ agrees with evolution under $H$ at time $t$.%
\footnote{\leanid{QuantumCircuit.heisenbergEvolution_inv_smul_mul}
    in \doclink{TNLean/Circuit/LiebRobinson/TimeScaling}, and
    \leanid{TNLean.MPS.AreaLaw.lemmaOneLength_eq_sourceLength}
    in \doclink{TNLean/MPS/AreaLaw/Hastings/LiebRobinson}.}
The same argument gives an admissible physical pair
$(10\Delta E\,u_0,8)$ for the enlarged length
$\tilde\xi=20u_0=\tfrac54\hat\xi$.%
\footnote{\leanid{TNLean.MPS.AreaLaw.ChainHamiltonian.norm_heisenberg_commutator_le_hastings_rescaled}
    and \leanid{TNLean.MPS.AreaLaw.lemmaOneEndpointLength_eq_sourceLength}.}
Since $6\tilde\xi=\tfrac{15}2\hat\xi$, the polynomial bound at
$\kappa=15/2$ gives
\begin{align}
    \|O_BO_LO_R-P_0\|\le K\tilde\xi^3e^{-l/(6\tilde\xi)},
    \label{eq:hastings07_lemma1_decay_length_endpoint}
\end{align}
with $K$ depending only on $R$. The endpoint is reached at
$\tilde\xi$, rather than at the original fixed length $\hat\xi$.%
\footnote{\leanid{TNLean.MPS.AreaLaw.exists_projectorFactorization_endpoint}
    in \doclink{TNLean/MPS/AreaLaw/Hastings/LemmaOneConstant}.}
At fixed $\hat\xi$, the factor $\sqrt{l/\hat\xi}$ in the appendix still
requires a strict increase of the decay length; the enlarged pair does
not invalidate that observation.

There is also a pair with velocity proportional to $J$ independently of
the gap. For $J>0$, take $v_*=40c_RJ$ and $\xi_{C,*}=80$. Its source
length is
\begin{align}
    s=\max(2v_*/\Delta E,80)
      =80\max(c_RJ/\Delta E,1)\ge2\tilde\xi.
    \label{eq:hastings07_lemma1_decay_length_source_enlarged}
\end{align}
Indeed, the maximum is at least each of $80c_RJ/\Delta E$ and $80$,
so it is at least their average $40(1+c_RJ/\Delta E)=2\tilde\xi$.%
\footnote{\leanid{TNLean.MPS.AreaLaw.two_mul_lemmaOneEndpointLength_le_sourceLength}.}
This pair allows the entropy argument to use an enlarged source length
while retaining a velocity of order $J$ and a numerical length scale.
The auxiliary velocities $8\Delta E u_0$ and $10\Delta E u_0$ contain
terms proportional to the gap. Although their time-cone estimates are
valid, their formulas are not fixed multiples of $J$. The estimates at
$\hat\xi$ and $\tilde\xi$ therefore do not directly establish the source's
stated velocity dependence; a separate bound relating the gap to $J$
would be needed to use these pairs for that purpose.

At the length $s$ in
(\ref{eq:hastings07_lemma1_decay_length_source_enlarged}), every conclusion
can use the same pair $(40c_RJ,80)$. Both contraction and positive
semidefinite factorizations have error at most
$Ks^3e^{-l/(6s)}$. The bootstrap inequalities, their Rényi forms, and the
entropy gain follow at $\xi'=6s$ with $C_1=Ks^3$, after increasing $K$
when necessary. For the matrix product state approximation, the entropy
maximum is the printed bound
$c_0\xi'\ln\xi'\ln q\,2^{\xi'\ln q}$, and the common integer threshold is
$\lceil e^{2S_{\max}}/2\rceil$. Thus the same source parameters enter the
factorization, entropy bound and approximation.
\footnote{The five source-parameter consequences are proved in
    \doclink{TNLean/MPS/AreaLaw/Hastings/SourceConsequences}.}


Replacing the window factor by $O_B^\dagger O_B$ gives a positive
semidefinite contraction at half the exponential rate. This is the
separate positivity correction in
\path{hastings07_lemma1_positivity_rate.tex}; the numerical propagation
bound does not remove it.

\section{The polynomial bound}

Lemma~1 requires $C_1(\xi)$ to be bounded by a polynomial in $\xi$.
\footnote{\path{References/0705.2024/law16.tex}, lines~258--262.}
In the unscaled statement the explicit coefficients are not homogeneous in
the energy scale: they contain terms such as $2/\Delta E$ and
$J^2/\Delta E$, so no bound in $\xi$ alone is possible there. In the
statement for $H/\Delta E$ every coefficient is a polynomial in
$x=J/\Delta E$, and $x\le\hat\xi$ with $\hat\xi\ge16$. The absorption of
$\sqrt l$ costs
$1/\sqrt{2e(1/(6\hat\xi)-1/(\kappa\hat\xi))}\le\hat\xi/\sqrt\beta$ with
$\beta=2e(\kappa-6)/(6\kappa)$, for $\xi'=\kappa\hat\xi$ and a fixed
$\kappa>6$. The short-length constant $e^{(12R+8)/(6\hat\xi)}$ is bounded.
Together these give
\begin{align}
    C_1\le K\hat\xi^3,
    \label{eq:hastings07_lemma1_decay_length_polynomial}
\end{align}
with $K$ depending only on $R$ and $\kappa$.\footnote{\leanid{TNLean.MPS.AreaLaw.lemmaOneConst_le_pow}
    and \leanid{TNLean.MPS.AreaLaw.exists_projectorFactorization_pow} in
    \doclink{TNLean/MPS/AreaLaw/Hastings/LemmaOneConstant}.}
For the source's nearest-neighbour terms ($R=1$) the constant $K$ depends
only on $\kappa$. The bound holds at the rate $\xi'=\kappa\hat\xi$ for
every $\kappa>6$, not at $\kappa=6$; the printed rate is reached at the
inflated length $\tilde\xi$ of (\ref{eq:hastings07_lemma1_decay_length_endpoint}).

\section{Conclusion}

The numerical prefactor, arbitrary second support, time-cone estimate,
physical-time rescaling and identification of the factorization lengths
are proved. The former quantitative restriction from \ghissue{9811} is
therefore removed. The reading of the printed Lieb--Robinson formula is
an explicit local correction. The Gaussian parameter correction, the
strict-rate issue at a fixed length and the positivity rate remain
recorded separately; no assertion is made that the literal printed
positive factorization follows at its original fixed length.

\bibliographystyle{alpha}
\bibliography{references}
\end{document}
